\section{Experimental Setup}
\label{sec:experimental-setup}

\paragraph{Datasets and retrieval.}
We evaluate on eight passage-ranking benchmarks. The TREC Deep Learning
datasets are TREC DL 2019 and TREC DL 2020 over the MS MARCO passage
corpus~\citep{bajaj2016msmarco,craswell2020dl19,craswell2021dl20}. To
test zero-shot transfer beyond MS MARCO-style web queries, we also use
six BEIR domains~\citep{thakur2021beir}: \textsc{DBPedia-Entity},
\textsc{NFCorpus}, \textsc{SciFact}, \textsc{TREC-COVID},
\textsc{Webis-Touche2020}, and \textsc{FiQA}. These datasets vary in
query style, corpus genre, domain specificity, and relevance density.

For every dataset, the first-stage run is BM25~\citep{robertson2009probabilistic}
from Pyserini/Anserini~\citep{lin2021pyserini,yang2017anserini}. We
rerank only the BM25 top-$N$ pool, with
$N \in \{10,20,30,40,50\}$. All methods receive the same first-stage
pool for a given query and pool size. For MaxContext, this means the
initial prompt window is $W=N$ and the output depth is $k=N$; the whole
BM25 top-$N$ pool is visible in the first call.

\paragraph{Models.}
The required model matrix contains nine open-weight causal LLMs from
three families: five Qwen3.5 models, one Llama-3.1 model, and three
Ministral-3 models. Specifically, we use
Qwen3.5-\{0.8B, 2B, 4B, 9B, 27B\}, Llama-3.1-8B-Instruct, and
Ministral-3-\{3B, 8B, 14B\}-Instruct-2512
~\citep{qwen35modelcard,dubey2024llama3,ministral3modelcard}. This
multi-family design is intended to test whether MaxContext is a
model-family-general reranking pattern rather than a behavior specific
to a single tokenizer, chat template, or long-context implementation.

All experiments use the instruction or post-trained chat variant of each
checkpoint and generation-based prompting with greedy decoding. We do not
use likelihood scoring for MaxContext, because the joint best--worst
objective in \textsc{MaxContext-DualEnd} is not equivalent to an
independent single-label likelihood score. When a model interface exposes
reasoning or thinking traces, we disable them so that outputs remain
short label selections rather than chain-of-thought generations.

\paragraph{Context-fit and inference invariants.}
Each candidate passage is rendered with passage length $\ell=512$. At
$N=100$, this gives approximately 26K input tokens before output reserve,
although the exact count is tokenizer-dependent. The published context
budgets of the selected model families are above this regime, but we do
not rely only on advertised limits: every rendered prompt is tokenized
before inference and must pass a context-fit preflight check. Prompts are
never truncated. A model--query--pool cell that fails this check is
blocked rather than silently shortened.

All MaxContext runs enforce the invariants from
\S\ref{sec:methodology}: numeric one-based document labels, strict
integer parsing, no fallback repair for invalid labels, no permutation
self-consistency, and deterministic BM25 resolution for the final
two-document residual in the single-extreme variants. Standard setwise
baselines use the same model, decoding configuration, first-stage pool,
and passage rendering, differing only in the local-window setwise
scheduling and sorting procedure.

\paragraph{Method grid.}
We evaluate seven reranking methods. Four are standard setwise baselines
following the local-window formulation of \citet{zhuang2024setwise}:
TopDown setwise with bubblesort, TopDown setwise with heapsort,
BottomUp setwise with bubblesort, and BottomUp setwise with heapsort.
The remaining three are the MaxContext variants introduced in this
paper: \textsc{MaxContext-TopDown}, \textsc{MaxContext-BottomUp}, and
\textsc{MaxContext-DualEnd}. Standard DualEnd sorting baselines are not
included in the EMNLP method matrix.

The main experiment is therefore:
\[
7\
\times 9\
\times 8\
\times 5\
= 2520\ \text{runs}.
\]
The primary efficiency contrast within the MaxContext family is at
$N=100$, where \textsc{MaxContext-DualEnd} uses $50$ LLM calls and either
single-extreme MaxContext control uses $98$ LLM calls plus a deterministic BM25 two-document endgame.

\paragraph{Evaluation metric.}
The primary effectiveness metric is per-query nDCG@10
~\citep{jarvelin2002cumulated}. We compute nDCG@10 after reranking the
BM25 top-$N$ pool and report macro averages over queries. We also retain
per-query scores for paired significance testing, confidence intervals,
and stability analysis. Because MaxContext produces a full permutation
of the BM25 pool, the same evaluation code is used for all seven
methods.

\paragraph{Statistical protocol.}
For each model--dataset--pool-size cell, we use paired bootstrap tests
over per-query nDCG@10 scores~\citep{efron1993bootstrap,
smucker2007comparison,urbano2019significance}. The bootstrap unit is the
query: each resample draws queries with replacement and recomputes the
mean paired difference between two methods. We report paired-bootstrap
95\% confidence intervals for method differences and use two-sided
paired-bootstrap $p$-values for significance annotations.

Multiplicity is controlled with Bonferroni correction across the
seven-method family within each model--dataset--pool-size cell
~\citep{bonferroni1936teoria}. This keeps the main significance protocol
simple and conservative for a short paper. Appendix tables report the
uncorrected paired differences, confidence intervals, and corrected
$p$-values so that the effect sizes remain visible even when a contrast
does not pass the corrected threshold.

\paragraph{Input-order stability.}
Whole-pool prompting may introduce position sensitivity, so we evaluate
input-order stability explicitly rather than assuming order invariance.
The required stability matrix uses DL19 with three representative models:
Qwen3.5-9B, Llama-3.1-8B-Instruct, and
Ministral-3-8B-Instruct-2512. For each of the seven methods and each
pool size $N \in \{10,20,30,40,50,100\}$, we run ten randomized input-order
repetitions using the same underlying BM25 pool. This gives
\[
7\
\times 3\
\times 1\
\times 5\
\times 10\
= 1050\ \text{stability runs}.
\]
We report mean nDCG@10, standard deviation, range, and worst-pair
differences across repetitions. We also analyze selected-label position
patterns for one representative model from each family to diagnose
whether whole-pool visibility changes positional behavior.

\paragraph{Cost and logging.}
For every run, we log LLM calls, prompt tokens, completion tokens, total
tokens, parser outcomes, BM25 bypass counts, per-query wall-clock time,
and aggregate wall-clock time. These axes are reported separately because
they measure different costs: reducing serial LLM calls does not imply
reducing input tokens, output tokens, FLOPs, or wall-clock time in every
setting~\citep{peng2025flopsreranking}. Our main efficiency claim is
therefore restricted to deterministic LLM-call counts and observed
wall-clock behavior.

All outputs are written as query-level JSONL logs and TREC-format run
files. The release package includes code, launchers, prompts, evaluation
scripts, logs, and a revision-pinned model manifest so that the exact
checkpoint, tokenizer, decoding, and prompt-rendering configuration can
be audited.